%% This template can be used to write a paper for
%% Computer Physics Communications using LaTeX.
%% For authors who want to write a computer program description,
%% an example Program Summary is included that only has to be
%% completed and which will give the correct layout in the
%% preprint and the journal.
%% The `elsarticle' style is used and more information on this style
%% can be found at 
%% http://www.elsevier.com/wps/find/authorsview.authors/elsarticle.
%%
%%
\documentclass[preprint,12pt]{elsarticle}

%% Use the option review to obtain double line spacing
%% \documentclass[preprint,review,12pt]{elsarticle}

%% Use the options 1p,twocolumn; 3p; 3p,twocolumn; 5p; or 5p,twocolumn
%% for a journal layout:
%% \documentclass[final,1p,times]{elsarticle}
%% \documentclass[final,1p,times,twocolumn]{elsarticle}
%% \documentclass[final,3p,times]{elsarticle}
%% \documentclass[final,3p,times,twocolumn]{elsarticle}
%% \documentclass[final,5p,times]{elsarticle}
%% \documentclass[final,5p,times,twocolumn]{elsarticle}

\usepackage{graphicx}
\usepackage{amssymb}
\usepackage{minted}
\usepackage{xcolor}
\definecolor{nicered}{rgb}{0.5,0.,0.}
\definecolor{nicegreen}{rgb}{0.,0.5,0.}
\definecolor{niceblue}{rgb}{0.,0.,0.5}
\usepackage{hyperref}
\hypersetup{colorlinks,citecolor=nicegreen,linkcolor=nicered,urlcolor=niceblue}
\usepackage{siunitx}

\newcommand{\dd}{\mathrm{d}}

%% The lineno packages adds line numbers. Start line numbering with
%% \begin{linenumbers}, end it with \end{linenumbers}. Or switch it on
%% for the whole article with \linenumbers after \end{frontmatter}.
%% \usepackage{lineno}

%% natbib.sty is loaded by default. However, natbib options can be
%% provided with \biboptions{...} command. Following options are
%% valid:

%%   round  -  round parentheses are used (default)
%%   square -  square brackets are used   [option]
%%   curly  -  curly braces are used      {option}
%%   angle  -  angle brackets are used    <option>
%%   semicolon  -  multiple citations separated by semi-colon
%%   colon  - same as semicolon, an earlier confusion
%%   comma  -  separated by comma
%%   numbers-  selects numerical citations
%%   super  -  numerical citations as superscripts
%%   sort   -  sorts multiple citations according to order in ref. list
%%   sort&compress   -  like sort, but also compresses numerical citations
%%   compress - compresses without sorting
%%
%% \biboptions{comma,round}

% \biboptions{}

%% This list environment is used for the references in the
%% Program Summary
%%
\newcounter{bla}
\newenvironment{refnummer}{%
\list{[\arabic{bla}]}%
{\usecounter{bla}%
 \setlength{\itemindent}{0pt}%
 \setlength{\topsep}{0pt}%
 \setlength{\itemsep}{0pt}%
 \setlength{\labelsep}{2pt}%
 \setlength{\listparindent}{0pt}%
 \settowidth{\labelwidth}{[9]}%
 \setlength{\leftmargin}{\labelwidth}%
 \addtolength{\leftmargin}{\labelsep}%
 \setlength{\rightmargin}{0pt}}}
 {\endlist}

\journal{Computer Physics Communications}

\begin{document}

\begin{frontmatter}

%% Title, authors and addresses

%% use the tnoteref command within \title for footnotes;
%% use the tnotetext command for the associated footnote;
%% use the fnref command within \author or \address for footnotes;
%% use the fntext command for the associated footnote;
%% use the corref command within \author for corresponding author footnotes;
%% use the cortext command for the associated footnote;
%% use the ead command for the email address,
%% and the form \ead[url] for the home page:
%%
%% \title{Title\tnoteref{label1}}
%% \tnotetext[label1]{}
%% \author{Name\corref{cor1}\fnref{label2}}
%% \ead{email address}
%% \ead[url]{home page}
%% \fntext[label2]{}
%% \cortext[cor1]{}
%% \address{Address\fnref{label3}}
%% \fntext[label3]{}

\title{\texttt{Candia-v2}: Logarithmic expansions for DGLAP evolution in $x$-space}

%% use optional labels to link authors explicitly to addresses:
%% \author[label1,label2]{<author name>}
%% \address[label1]{<address>}
%% \address[label2]{<address>}

\author[a]{First Author\corref{author}}
\author[a,b]{Second Author}
\author[b]{Third Author}

\cortext[author] {Corresponding author.\\\textit{E-mail address:} firstAuthor@somewhere.edu}
\address[a]{First Address}
\address[b]{Second Address}

\begin{abstract}
%% Text of abstract
A submitted program is expected to satisfy the following criteria: it must be of benefit to other physicists, or be an exemplar of good programming practice, or illustrate new or novel programming techniques which are of importance to computational physics community; it should be implemented in a language and executable on hardware that is widely available and well documented; it should meet accepted standards for scientific programming; it should be adequately documented and, where appropriate, supplied with a separate User Manual, which together with the manuscript should make clear the structure, functionality, installation, and operation of the program.

Your manuscript and figure sources should be submitted through Editorial Manager (EM) by using the online submission tool at \\
https://www.editorialmanager.com/comphy/.

In addition to the manuscript you must supply: the program source code; a README file giving the names and a brief description of the files/directory structure that make up the package and clear instructions on the installation and execution of the program; sample input and output data for at least one comprehensive test run; and, where appropriate, a user manual.

A compressed archive program file or files, containing these items, should be uploaded at the "Attach Files" stage of the EM submission.

For files larger than 1Gb, if difficulties are encountered during upload the author should contact the Technical Editor at cpc.mendeley@gmail.com.
\\

% All CPiP articles must contain the following
% PROGRAM SUMMARY.

\noindent \textbf{PROGRAM SUMMARY/NEW VERSION PROGRAM SUMMARY}
  %Delete as appropriate.

\begin{small}
\noindent
{\em Program Title:} Candia-v2                       \\
{\em CPC Library link to program files:} (to be added by Technical Editor) \\
{\em Developer's repository link:} https://github.com/champso1/candia-v2 \\
{\em Licensing provisions(please choose one):} GPLv3  \\
{\em Programming language:} C++/Fortran                        \\
{\em Supplementary material:} \textcolor{red}{TODO}    \\
  % Fill in if necessary, otherwise leave out.
{\em Journal reference of previous version:} https://doi.org/10.48550/arXiv.2512.22667    \\
  %Only required for a New Version summary, otherwise leave out.
{\em Does the new version supersede the previous version?:} Yes \\
  %Only required for a New Version summary, otherwise leave out.
{\em Reasons for the new version:} Availability of N3LO operator matrix elements, matching conditions, and approximate kernels \\
  %Only required for a New Version summary, otherwise leave out.
{\em Summary of revisions:} Conversion from C to C++, optimizations with convolution/interpolation routines, N$^3$LO evolution, and possibility for LHAPDF inputs \\
  %Only required for a New Version summary, otherwise leave out.
{\em Nature of problem(approx. 50-250 words):} This program solves the DGLAP evolution equations in $x$-space for parton distributions functions up to approximate N$^3$LO \\
  %Describe the nature of the problem here. \\
{\em Solution method(approx. 50-250 words):} The algorithm is based upon a logarithmic expansion of the solution in $x$-space and a set of recursion relations for unknown coefficients. \\
  %Describe the method solution here.
{\em Additional comments including restrictions and unusual features (approx. 50-250 words):}\\
  %Provide any additional comments here.
  Due to the dependency on the GNU Scientific Library (GSL), which doesn't provide native binaries for Windows and is also difficult to build from source, WSL is recommended if using on Windows.
   \\
   


\begin{thebibliography}{0}
\bibitem{1}Reference 1         % This list should only contain those items referenced in the                 
\bibitem{2}Reference 2         % Program Summary section.   
\bibitem{3}Reference 3         % Type references in text as [1], [2], etc.
                               % This list is different from the bibliography at the end of 
                               % the Long Write-Up.
\end{thebibliography}
* Items marked with an asterisk are only required for new versions
of programs previously published in the CPC Program Library.\\
\end{small}
\end{abstract}
\end{frontmatter}

\section{Introduction}
\label{sec:intro}

We present a generalization of the $x$-space \texttt{Candia} algorithm to next-to-next-to-next-to-leading order (N$^3$LO) accuracy in Quantum Chromodynamics (QCD) for solving the DGLAP evolution equations for unpolarized parton densities in the nucleon.
The algorithm is based on logarithmic expansions of the solution and can be extended to all orders in QCD. 
An expansion equivalent to the exact solution of the DGLAP equation at N$^3$LO is presented in the non-singlet sector. 
Results for approximate N$^3$LO PDFs, evolved using the most recent approximations to the N$^3$LO DGLAP splitting functions, 
are provided for benchmarking. The new version of the code, \texttt{Candia-v2}, 
is publicly available at \url{https://github.com/champso1/candia-v2}~\cite{candiav2}.

\section{Description of the Program}
\label{sec:prog-desc}

\subsection{Requirements and Dependencies}
\label{subsec:reqs-and-deps}

There are some version requirements and dependencies one needs to install in order to successfully compile and run the code. They are:
%
\begin{itemize}
  \item \href{https://cmake.org/}{CMake} >= 3.20,
  \item C++20 compatibility. Among the standard compilers, one would need GCC >= 13 or CLang >= 17. Other compilers like Apple CLang or Intel compilers haven't been tested; MSVC is incompatible and will not work,
  \item A Fortran compiler such as GFortran,
  \item GSL (GNU Scientific Library): Required as a dependency of the \texttt{libome} package,
  \item (Optional) LaTeX utilities (PDFLatex/Biber): if one wants to build this manuscript and other relevant ones included in the repository,
  \item (Optional) \href{https://www.doxygen.nl/}{Doxygen}: if one wants to build the code documentation.
  \item (Optional) \href{https://lhapdf.hepforge.org/}{LHAPDF}: if one wants to interface with LHAPDF
\end{itemize}

For the latter three optional dependencies, see Section~\ref{subsec:acq-and-build} for descriptions of the CMake flags needed to enable these features.

\subsection{Acquiring and Building}
\label{subsec:acq-and-build}

The code is publically available on \href{GitHub}{https://github.com/champso1/candia-v2}.
There are a number of active branches in which various new features are being added/tested, but by default the \mintinline{cfg}{main} branch contains the most stable version. While one can at any time build from this branch, it is preferrable to either download the source from the Releases page, or checkout the corresponding tag from the cloned repository.
After acquiring the source code, one follows the standard CMake building process:
%
\begin{minted}{bash}
git clone https://github.com/champso1/candia-v2
cd candia-v2
git submodule init 
git submodule update
mkdir build && cd build
cmake .. [-G <Generator>] \
         [<other options>]
cmake --build .
cmake --install .
\end{minted}
%
Note the commands related to git submodules, which will correctly initialize the \texttt{libome} repository.
Among the optional \mintinline{bash}{<other options>} are the standard CMake options such as \mintinline{cfg}{CMAKE_CXX_COMPILER} or \mintinline{cfg}{CMAKE_BUILD_TYPE} if one wants to build in anything other than Release mode, which is the default.
There are also a few additional options that are specific to this program. They are:
%
\begin{itemize}
\item \textbf{\mintinline{cmake}{ENABLE_THREADING}}: Due to the nature of the algorithm, the non-singlet distributions are entirely independent, and in the singlet sector, the grid can be split up and evolved independently as well. This option enables building the code with the appropriate threading libraries, allowing the program to place these independent evolutions on separate threads such that they run concurrently.
  This option is off by default for compatibility with older systems, but is recommended to be enabled if the computer has multiple CPU cores.
\item \textbf{\mintinline{cmake}{BUILD_MANUSCRIPTS}}: Contained in the \mintinline{cfg}{manuscripts} directory in the root of the project is this manuscript as well as the manuscript submitted to Physical Review D, containing a higher level description of the algorithm. If this option is enabled, then a new target \mintinline{cfg}{manuscripts_pdf} is generated which will build both of these manuscripts.
  Requires LaTeX utilities (PdfLaTeX, Biber) and some standard packages.
  This option is off by default.
\item \textbf{\mintinline{cmake}{BUILD_DOCS}}: Use Doxygen to build the documentation for the code. This option is off by default.
\item \textbf{\mintinline{cmake}{WITH_LHAPDF} / \mintinline{cmake}{LHAPDF_ROOT_DIR}}: Enabling the former will compile in all the features relying on LHAPDF; the latter is optional unless if the user has installed LHAPDF into a nonstandard location, that being anywhere but \mintinline{bash}{/usr/local}, the default installation directory. These options are off/empty by default.
\end{itemize}

\subsection{Usage: Example Program and Building/Linking}
\label{subsec:ex-prog}

In Listing~\ref{lst:ex-prog} is a full C++ program that will perform the DGLAP evolution at N$^3$LO starting from the standard benchmark initial conditions.
The output is stored in an \mintinline{cpp}{std::vector<ArrayGrid>}, which the user can handle however they need; the \mintinline{cpp}{ArrayGrid} is a class that, for the purposes of accessing elements, behaves like an ordinary array, e.g. it has the \mintinline{cpp}{operator[]} method. See \textcolor{red}{TODO} for a further description of this class.
%
\begin{listing}
  \inputminted[frame=lines,linenos]{cpp}{res/example.cpp}
  \caption{Example program that performs a full DGLAP evolution at N$^3$LO starting from the standard benchmark initial conditions and ending at $\qty{100}{\GeV}$.}
  \label{lst:ex-prog}
\end{listing}
%
In the code, the user first includes the header \mintinline{cpp}{Candia-v2/Candia.hpp}. This is the main header to include, and all required utilities the user need interact with is included within that header.
We also define the following variables:
\begin{itemize}
  \item \mintinline{cpp}{num_grid_points}: the number of points to place on the interpolation/convolution grid.
  \item \mintinline{cpp}{order}: the perturbative order of the solution. this number is the $m$ in N$^m$LO.
  \item \mintinline{cpp}{Qf}: the final energy to evolve to.
  \item \mintinline{cpp}{kr}: the ratio of factorization and renormalization scales.
  \item \mintinline{cpp}{iterations}: the number of total iterations to complete. this number corresponds to the $s$ index \textcolor{red}{TODO}
  \item \mintinline{cpp}{trunc_idx}: the truncation index to use for the singlet solution. this number corresponds to the variable $\kappa$ \textcolor{red}{TODO}.
\end{itemize}
%
We also define an \mintinline{cpp}{std::vector} called \mintinline{cpp}{xtab}, which specifies a set of points which are directly placed on the grid and are able to be directly retrieved later, for benchmarking purposes. This avoids interpolating to get the value of the distributions at these points. Additionally, depending on the option passed to the \mintinline{cpp}{Grid} constructor, it will also guide the creation of the logarithmic/linear bins on the grid.
We then create the \mintinline{cpp}{Distribution} which contains the initial value of $\alpha_s$, the initial energy, as well as the functions defining the initial values for the distributions. \mintinline{cpp}{LesHouchesDistribution} defines the set of initial conditions used in \textcolor{red}{TODO}.
We use these values to create an \mintinline{cpp}{AlphaS} object, containing the information to evolve/match $\alpha_s$.
Lastly, we create a \mintinline{cpp}{DGLAPSolver} object, and pass it these parameters, and call \mintinline{cpp}{.evolve()} to perform the evolution and grab the distributions.

\subsection{Program Structure}
\label{subsec:prog-struct}

We now describe the main implementation details of the source code itself.

\subsubsection{The \texorpdfstring{\mintinline{cpp}{ArrayGrid}}{ArrayGrid} Class}
\label{subsec:arraygrid}

The \mintinline{cpp}{ArrayGrid} class is simply a wrapper over an ordinary \mintinline{cpp}{std::vector} with a few convenience methods, though largely behaves identically. This means that one can do, for instance, list-initialization and access using the \mintinline{cpp}{operator[]} method, as seen in Listing~\ref{lst:arraygrid-ex}.
%
\begin{listing}[h!]
\begin{minted}{cpp}
ArrayGrid arr1{1.0, 2.0, 3.0}; // list-initialization
ArrayGrid arr2(20); // 20 elements, all zeroed
arr2[0] = arr1[0]; // const and non-const ref
for (double& x : arr1) // begin/end iterators
    x *= x;
\end{minted}
\caption{Demonstration of the \mintinline{cpp}{ArrayGrid} class.}
\label{lst:arraygrid-ex}
\end{listing}
%
This is used as the datatype for the PDFs.

\subsubsection{Splitting Functions and OMEs}
\label{subsec:splitfuncs-and-omes}

Structurally, the splitting functions and the OMEs have the following structure:
%
\begin{equation}
  A(x) + \frac{B(x)}{(1-x)^+} + C\delta(1-x),
\end{equation}
%
i.e. a "regular" piece, given by the $A$ coefficient, a "plus" piece, given by the $B$ coefficient (\textit{not} including the plus-distribution itself), and a "delta" piece, given by the $C$ coefficient (also \textit{not} including the delta function).
For the purposes of convolutions with these objects, we can keep one generic convolution function and utilize C++ OOP to define a generic \mintinline{cpp}{Expression} class that contains these three methods.
We then have that the classes \mintinline{cpp}{SplittingFunction} and \mintinline{cpp}{OpMatElem} inherit from this and implement the required 
%
Importantly, since the convolutions proceed via a Gauss-Legendre quadrature, there are a set of fixed points at which these objects are evaluated, which is known before the evolution begins. Therefore, \mintinline{cpp}{Expression} objects contain a cache (implemented as an \mintinline{cpp}{std::map}) of the expression evaluated at those points that is filled at the start of the evolution. This is especially helpful for higher orders where the expressions become very complex, and it saves a large amount of time.

\subsubsection{Initial Distributions}
\label{subsec:init-dists}

Initial distributions supply the values of the parton distributions at some initial energy $Q_0$, as well as the initial value of $\alpha_s$. Since this behavior is independent of the evolution, it is factored into a \mintinline{cpp}{Distribution} class.
Obviously, there are many ways to provide initial conditions, so it is a virtual class that specifies a set of behavior that is supplied by the user's chosen conditions:
the values of the quark masses, the initial value of $\alpha_s$, functions for $xf_p(x)$ where $p$ is any parton, the initial value of $n_f$, the number of massless flavors at $Q_0$, as well as $Q_0$ itself.
%
We have provided two derived classes, \mintinline{cpp}{LesHouchesDistribution} and \mintinline{cpp}{LHAPDFDistribution}.
The former implements the conditions used in the benchmarks \textcolor{red}{TODO}, while the former provides an interface to supply initial conditions from an LHAPDF PDF.
Note that the latter option is only present in the compiled version of the code if the user enables the \mintinline{cmake}{WITH_LHAPDF} and, if necessary, the \mintinline{cmake}{LHAPDF_ROOT_DIR} CMake variables during configuration.

\subsubsection{The \texorpdfstring{\mintinline{cpp}{AlphaS}}{AlphaS} Class}
\label{subsec:alphas}

The behavior of the QCD running coupling $\alpha_s$ is stored within the \mintinline{cpp}{AlphaS} class, including the values of the various $\beta$-coefficients and a routine to evaluate the coupling at some other arbitrary scale given its value at an initial scale. This is done by solving the Cauchy problem
%
\begin{equation}
  \frac{\dd \alpha_s(t)}{\dd t} = \beta(\alpha_s).
\end{equation}
%
At LO, this has an exact solution, while at higher orders a fourth-order Runge Kutta method is used.
Additionally, this class also contains the necessary information and methods to perform the matching of the coupling at the quark thresholds at NNLO and above, assuming the on-shell scheme.
It keeps an internal set of arrays that contain the values of the coupling pre- and post-threshold for each relevant $n_f$; for this reason, its constructor requires the mass array which is suppied by the initial distribution.

\subsubsection{The Grid}
\label{subsec:grid}

The interpolation/convolution grid is arguably the most important part of the code, as convolutions and interpolations take up the vast majority of the program's runtime. At the highest level, the \mintinline{cpp}{Grid} class supplies an underlying grid of points, along with methods to perform convolutions between arrays and expressions as well as interpolations of arrays on the grid. These have the signature
%
\mint{cpp}|double Grid::convolution(ArrayGrid& A, Expression& E, uint k)|
%
\mintinline{cpp}{A} is an array that represents an ordinary function that is assumed to be defined at the same points as the grid; \mintinline{cpp}{E} represents a splitting function or operator matrix element, and $k$ is the grid index at which to perform the convolution.
%
\mint{cpp}|double Grid::interpolate(ArrayGrid& A, double x)|
%
\mintinline{cpp}{A} is an array like above, and $x$ is the point one wants to interpolate.
%
There are several ways to define the grid.
The original version of \texttt{Candia} used logarithmically spaced intervals, extending from the minimum value provided in the array of tabulated $x$-values to 1. While this works well due to prescription we apply to do the convolutions, it leaves much fewer points for $x>0.5$, which is a particularly sensitive region especially at higher orders where PDFs become very small.
Because of this, we have devised a few other prescriptions in an attempt to fill the higher end of the grid with more points and increase the accuracy in this regime.
The first will create several more logarithmic bins for $x>0.1$, in an attempt to fill this region with more points. The other prescription will also still retain the ordinary logarithmic bins for $x<0.1$, but for $x>=1.0$ it will linearly space the points.

\subsubsection{The Evolution}
\label{subsec:evol}

The part of the code responsible for the actual evolution of the parton distributions ties all of the above classes and methods together; this is all done within the \mintinline{cpp}{DGLAPSolver} class. The constructor accepts the following values:
\begin{itemize}
\item the perturbative order
\item a \mintinline{cpp}{Grid} object
\item a \mintinline{cpp}{AlphaS} object
\item $Q_f$, the final evolution energy ($Q_0$ is encoded in the \mintinline{cpp}{AlphaS} object)
\item the number iterations (corresponding to the $s$ index)
\item the number of terms to keep in the singlet expansion, called the \textit{truncation index} (this is $\kappa$)
\item a \mintinline{cpp}{Distribution} object
\item the ratio $\mu_r^2/\mu_f^2$ (this defaults to 1 if the user doesn't provide an argument)
\end{itemize}
%
Upon construction, it will store copies of the important arguments, then allocate memory for and fill the initial part of the distributions from the \mintinline{cpp}{Distribution} object. At N$^3$LO, it will also hard-code the $r_1$, $b$, and $c$ coefficients.
Once the user calls the \mintinline{cpp}{.evolve()} method, the evolution proceeds. It will first form the special distributions the are evolved, e.g. $q_{NS}^{+}$ or $q^{(+)}$. It will then, using the initial value of $n_f$ from the distribution and the corresponding values of $\alpha_s$, perform the evolution up until the first relevant quark mass. This is accomplished by calculating the logarithmic coefficients, then using the recursion relations to calculate, up until the specified number of iterations, the final values of the coefficients, then perform the resummation. Then, it will reform the quark PDFs, and if relevent, it will perform the matching, before proceeding again.
If the final requested evolution energy is less than a quark mass, instead of proceeding all the way to the quark mass, it will proceed only to the requested value, skipping matching and storing the results of the resummation in a special set of distributions that are then returned to the user.

These steps are performed if the user requests to perform the evolution in the variable flavor number scheme. If the user decides to perform in the fixed flavor number scheme, we proceed directly from the initial to the final energy, calculating quantities related to $\alpha_s$ only with the requested number of flavors.

\subsection{Major Optimizations from \texorpdfstring{\texttt{Candia-v1}}{Candia-v1}}
\label{subsec:optims}

We now, for completeness, describe the major optimizations and general enhancements that have been made with the advance to \texttt{Candia-v2}. The first one is the flexibility offered by C++ over C when it comes to type-generic functions and interfaces. While this can be achieved in C either directly or indirectly, it often leads to un-transparent code that is hard to read and debug. C++ provides a very natural way to express, say, ``Expressions'', as in Section~\ref{subsec:splitfuncs-and-omes} with neglible overhead. While this doesn't have any direct performance impacts, it makes debugging and handling the corresponding code very straightforward, leading to a faster developing process.

There are three major improvements that were made that significantly increased the performance of the code, those being the caching of the values of the splitting functions and operator matrix elements, the usage of multi-threading due to the independent nature of certain pieces of the evolution, and the memory reduction by only allocating for two iterations, as opposed to all $s$ requested ones.

In regards to the first, the choice of Gauss-Legendre quadrature leaves a fairly manageable memory footprint if we simply cache the values of the splitting functions evaluated at the Gauss-Legendre abscissae.
Additionally, it seemed quite wasteful to perform these calculates hundreds if not thousands of times during the course of the program's execution.
For NNLO and beyond, this reduces several dozen calculations and calls to harmonic polylogarithm routines to a few conditional statements to retrieve the value from the cache, and drastically speeds up the calculation of the convolutions.

In regards to the second, the non-singlet distributions that are evolved are completely decoupled, meaning they can all proceed concurrently. Assuming the user enables the \mintinline{cmake}{ENABLE_THREADING} CMake option, the code is then able to place each individual non-singlet evolution on a separate thread. The singlet sector, unfortunately, doesn't have this freedom. However, each iteration requires information only from the previous iteration, not from it's own, when it comes to performing the convolutions, the most time-intensive part of the convolution. Therefore, we can split the grid up into several sections and have the convlutions performed independently on these parts, reducing the time spent on the convolutions for that iteration by, roughly, the number of threads. Since essentially all remotely recent CPUs have multiple cores/threads, this is very much a viable option for performance.

In regards to the third and final optimization, the realization about the singlet sector above is also present in the non-singlet sector, though only for the outermost index, denoted $s$. Previously, in \texttt{Candia-v1}, this was ignored, and memory was allocated for all iterations, meaning, for instance, at NNLO, if the user requested 10 iterations and 500 grid points, we would request $8*37*10^3*500 = 148000000$ bytes of memory (8 being the size of a double-precision floating point on most systems and 37 being the total number of distributions used throughout the evolution), which is approximately \qty{148}{\mega\byte}. With this optimization, we drop the memory usage by a factor of $s/2$, which is very impactful for large grids and large numbers of iterations.


\bibliographystyle{elsarticle-num}
\bibliography{sources}


\end{document}
